\section{\sysname Overview}
\label{sec:challenges}

\begin{figure}
    \centering
    \includegraphics[width=\linewidth]{figures/arch_figure/Tuft-TuFT arch.pdf}
    \caption{\sysname architecture in one model group.}
    \label{fig:arch}
    \Description{}
\end{figure}

Given these challenges, \sysname is organized into two components that work together (Fig.~\ref{fig:arch}): the hybrid backends (\S\ref{sec:backend}) and the RL-loop-aware scheduler (\S\ref{sec:sched}). The hybrid backends make the train/sample boundary a runtime variable, flipping a group between the two roles at zero copy while a small sampling-side Corrector keeps the returned log-probabilities faithful to each tenant's own policy. The scheduler infers each tenant's loop online and rearranges its sampling so that training bursts pack into contiguous blocks, opening long concentrated sampling windows and keeping the flips rare.

A request enters through the router and is directed to its model group. Training requests carry data, a loss function and an adapter id; loop discovery observes them to update its estimate of the tenant's period and burst length, which drives the scheduler's decisions. Sampling requests carry data and an adapter id, and are routed to the corresponding fixed-sampling group or flex-sampling group based on that decision. Whether a flex group is currently serving samples or running training is set by the scheduler, so the router sees a set of sampling groups whose composition changes over time.

\section{Hybrid Backends with High-fidelity Enabling Flexible GPU Flow}
\label{sec:backend}



\subsection{Two backend types and the asymmetry design}
\label{sec:twobackends}

Because of the imbalance nature of sampling and training and the memory waste caused by such imbalance, we proposed an asymmetry structure to perfectly solve such imbalance and resource waste. The general structure of reinforce learning infrastructure is to spawn disaggregated sampling and training backends. We notice the sampling's dominance, so we use a set of GPU to serve sampling only, called fixed sampling backend. We notice the potential training periodicity and training waste, so we use a set of GPUs to serve both training and sampling as flex backend, which can flip between the two status flex-training and flex-sampling.

Throughout, each backend occupies a \emph{group}: a set of GPUs holding one
tensor-parallel copy of the base---a TP group, i.e.\ a single model replica. A group is
the atomic unit that flips role and that requests are routed to; data parallelism, whether
across flex groups during training or across sampling groups, is an \emph{arrangement of
groups} rather than a unit of its own.

\paragraph{Fixed sampling backend.} This backend serves only for sampling. It can
also enable quantization when needed. When no quantization is enabled, the only
difference between the fixed sampling backend and a flex backend in flex-sampling is
\emph{who owns the weights}: the fixed backend loads and owns its own copy of the base
model, whereas a flex backend owns no weights at all---it aliases the copy held by the
training side (\S\ref{sec:zerocopy}). Two consequences follow, and together they are
what make the fixed backend worth its GPUs. First, a backend that never gives up its
role never gives up its KV cache either; its paged KV, prefix cache and adapter cache
stay warm indefinitely, so it is the stable sampling capacity that the scheduler can
always count on. Second, because it is not required to keep its weights bit-identical
to the trainer's, it is free to apply inference-only optimizations that a
weight-sharing group cannot. Quantization changes both the dtype and the layout of the
served weights and would immediately invalidate the alias, so it is available on the
fixed path and structurally unavailable on the flex path.

\paragraph{Flex backend.} This backend serves mainly for training. It can flip between training and sampling. The traing and sampling use the same base model but use different computation framework. It can achieve both the best training and sampling performance but the minimum switch overhead. With the RL loop scheduler, the training will be organized as tight training windows. The rest time will be used for sampling.

The asymmetry is deliberate and one-sided: there is a sampling-only backend but no
training-only backend. Sampling exceeds training by one to two orders of magnitude
(\S\ref{sec:charact:sampling_training_ratio}), so dedicating a group permanently to
sampling is always profitable, whereas dedicating one permanently to training would
strand capacity for most of every RL period. Every group that can train is therefore a
flex group, and \sysname keeps at least one of them training-capable at all times so
that a tenant's \textsf{update} never waits for a re-composition. Seen this way the two
backend types are not two scheduling states but two \emph{optimization regimes}: the
flex regime is constrained to the trainer's layout and numerics, which is the price of
a zero-copy flip, and the fixed regime is unconstrained, which is the reward for never
flipping.

\subsection{Zero-copy switching at a shared layout}
\label{sec:zerocopy}

\begin{figure}
    \centering
    \includegraphics[width=\linewidth]{figures/arch_figure/Tuft-GPU transfer.pdf}
    \caption{GPU transfer of fast loop and slow loop. The fast loop flips a flex group
    between flex-training and flex-sampling with no weight movement; the slow loop
    re-composes Fixed$\leftrightarrow$Flex, which reloads weights and may change their
    numeric format.}
    \label{fig:gputransfer}
    \Description{}
\end{figure}

A flex backend must present the same base model to two different engines: a PyTorch
tensor-parallel trainer and a vLLM sampler. The naive way to do this is to materialize
the weights twice---either as two resident copies, which halves the memory available
for KV cache, or as one copy that is handed over by tearing down one runtime and
standing up the other on every flip, which costs hundreds of seconds
(\S\ref{sec:GPU_waste}). \sysname does neither. It keeps exactly
\emph{one} copy of the base weights in device memory for the lifetime of the group and
lets whichever engine is active reference it in place.

\paragraph{Making the two layouts identical.} The enabling observation is the layout
coincidence of \S\ref{sec:tp-convergence}: at the same tensor-parallel degree, a
trainer shard and a vLLM shard hold the same slice of the same tensor. \sysname makes
this exact rather than approximate by fusing the trainer's linear layers the way vLLM
fuses them---$q,k,v$ into a single \texttt{qkv\_proj} and gate/up into a single
\texttt{gate\_up\_proj}---and packing the fused weights rank-major, so that the local
shard held by trainer rank $i$ is byte-for-byte the shard vLLM expects on
tensor-parallel rank $i$. Column-wise parallelism is applied to the fused projections
and row-wise parallelism to \texttt{o\_proj} and \texttt{down\_proj}. Because the base
is frozen throughout RL fine-tuning and only LoRA parameters carry gradients, this
shard is not merely identical at the first flip but stable across all of them.

\paragraph{Aliasing instead of copying.} The sampler is started with a \emph{dummy}
weight load: vLLM allocates weight-shaped buffers and completes its full initialization
path---KV cache profiling, graph capture---without ever reading a checkpoint. \sysname
then exports each trainer shard as a CUDA IPC descriptor and, inside the sampler
process, rebinds each vLLM parameter to the producer's physical pages. No bytes move,
no host round-trip occurs, and the invariant that only one copy of the base exists in
memory is preserved. Two details make this survive vLLM's own optimizations. The
injection must happen \emph{before} CUDA graph capture, because a captured graph records
device pointers and would otherwise replay against the abandoned dummy buffers; we
therefore inject from a worker hook that runs between KV cache initialization and graph
capture. And the displaced dummy allocations are re-tagged rather than freed, so the
allocator keeps tracking their virtual addresses. Because the base is frozen, the
descriptors themselves are reusable: every flip after the first skips descriptor
construction entirely.

\paragraph{Flipping.} With the alias in place, a flip never touches weights. To enter
training, \sysname releases the physical pages behind the KV cache and behind the
displaced dummy weights through vLLM's virtual-memory allocator, returning that memory
to the trainer for activations and optimizer state; the \emph{virtual} address records
and the captured CUDA graphs are retained. To return to sampling, the same handles are
re-mapped to freshly allocated physical pages at their original virtual addresses.
Because the addresses are unchanged the graphs remain valid and are never re-captured,
and because the pages carry no state that must survive---KV is rebuilt by inference and
the dummy weights are never read---the re-map transfers no data. Both page classes must
be re-mapped together: leaving the weight-tagged pages unmapped faults, because vLLM's
non-parameter buffers and attention workspaces still hold references into that pool. A
flip is therefore a page-table operation whose cost scales with the amount of memory
reclaimed rather than with model size.

\begin{table}[t]
\centering
\small
\caption{Measured flip cost on A100-80GB. ``Reclaimed'' is the device memory returned
to the trainer per rank. No weight bytes are copied in either direction.}
\label{tab:flipcost}
\begin{tabular}{lrrr}
\toprule
Model / TP & Reclaimed & To training & To sampling \\
\midrule
Qwen3-4B, TP=2  & 39.1\,GB & 96\,ms  & 29\,ms  \\
Qwen3-32B, TP=4 & 12.0\,GB & 144\,ms & 732\,ms \\
\bottomrule
\end{tabular}
\end{table}

Table~\ref{tab:flipcost} reports the measured cost. Both directions complete well under
a second at both scales, against the hundreds of seconds a process-level teardown
requires (\S\ref{sec:GPU_waste}). The 32B row also exposes the one knob the mechanism
offers: releasing more KV gives the trainer more room but makes the return trip slower,
so \sysname releases the smallest amount that satisfies the training step's memory
requirement. Retaining the CUDA graphs is what makes the flip worth performing at all.
Capturing them costs minutes, and running the sampler eagerly instead would forfeit
$3.1\times$ throughput (752 versus 2318 tokens/s on 4B at TP=2).

\paragraph{What the flip gives up.} Dropping KV is not free: sequences in flight when a
flip begins lose their cache. \sysname does not pay this in practice because the
scheduler stops admitting work to a group as its flip approaches, so the cache is
already nearly empty when the flip happens (\S\ref{sec:sched:affix}). This is the
concrete sense in which the two pillars depend on each other: the flip is cheap because
it is scheduled, and scheduling is feasible because the flip is cheap.

\subsection{Mismatch Corrector: delivering the fidelity SLO}
\label{sec:mismatch-pred}

Zero-copy switching guarantees that the trainer and the sampler read the \emph{same
bytes}. It does not guarantee that they compute the same numbers. Paged attention,
different accumulation orders, chunked prefill, graph-level fusion and---on the fixed
path---quantization all mean that the log-probability vLLM reports for a token differs
from the one the trainer computes for that same token under the same weights.
\S\ref{sec:mismatch-char} measured this gap: it is present at zero staleness, present
on a backend that never switches, and amplified $2.5\times$ by quantization. It is the
structural cost of running two engines, paid by every disaggregated RL system and not
only by ours.

\paragraph{Why a small per-token gap is a real problem.} Write $\Delta_t = s_t - p_t$
for the difference between the sampling and the training log-probability of response
token $t$. The per-token mean is tiny---$c \approx 0.0023$ in bf16 and $0.0052$ under
FP8---but policy-gradient methods do not consume tokens, they consume sequence
likelihood ratios. The importance weight a tenant computes is
$w = \exp(\sum_t \Delta_t) \approx \exp(T\!\cdot\!c)$, which grows with response length:
a 500-token response is up-weighted $3.1\times$ in bf16 and $13.4\times$ under FP8
purely because it is long. The bias is therefore not noise that averages out but a
systematic preference for long sequences, and because it is signed it also breaks the
symmetry of any two-sided trust region. The damage is algorithm-dependent: methods that
apply a sequence-level ratio---REINFORCE, RLOO, off-policy GRPO---are affected directly,
whereas strictly on-policy token-level clipping absorbs $c$ within its clip range.

\paragraph{What is learnable, and what is not.} The natural first attempt is to predict
$\Delta_t$ per token and subtract it. This fails, and the reason dictates the design.
Decomposing the variance of $\Delta_t$ over a held-out corpus of 465K tokens shows that
only $0.80\%$ of it lies \emph{between} sequences; $99.2\%$ is within-sequence variation
with negligible autocorrelation. Two oracle experiments confirm the consequence.
Replacing $\Delta_t$ by its per-sequence mean drives the importance-weight error to
exactly zero, and a per-token oracle does no better; conversely, subtracting the
per-sequence mean \emph{raises} token-level MAE by $16.7\%$. Per-token accuracy is thus
worthless for the quantity that matters, and the within-sequence residual is better left
in place---it is what makes a trust region fire. The learnable and harmful component is
exactly one scalar per sequence.

\paragraph{Design.} The Corrector is a small network that consumes only information
available on the sampling side---the response token ids and their positions, the
sampling log-probabilities, the prompt length, the temperature and the adapter
rank---and predicts $\widehat{\Delta}$; the correction actually applied is
its per-sequence mean $\hat{c}$, so the log-probability \emph{returned} to the tenant is
$s_t - \hat{c}$. It never alters the sampled tokens and it never touches the trainer.
Three properties of the target justify this shape. The mismatch is \emph{user-agnostic}:
it is a property of the engine pair and the serving configuration, not of who is
training, so the network takes no tenant or task identifier and one Corrector serves
every tenant on a base model. It is \emph{knob-composable}: the dependence on
temperature and adapter rank is captured by taking them as continuous inputs rather than
by fitting one model per configuration. And it is \emph{stable} over the course of a
run, so it can be profiled offline rather than retrained online. Deployment follows
directly: one Corrector per base model per sampling path, calibrated once from a short
probe of synchronous rollouts for which both log-probabilities are available. \sysname
has two such paths---the flex path, which shares the trainer's bf16 weights, and the
fixed path, which may be quantized---and they require separate calibration precisely
because quantization widens the gap.

The output head is zero-initialized, so an untrained Corrector reproduces the
uncorrected system bit-exactly and any measured gain is attributable to what it learned.
The training objective combines a per-token term with a length-normalized per-sequence
term and an explicit bias term, reflecting that the sequence-level bias rather than
token accuracy is the quantity to minimize.

\paragraph{What the Corrector buys, and what it does not.} Its purpose is correctness,
not efficiency. When the standard deviation of $\sum_t \Delta_t$ greatly exceeds the
trust-region threshold---the regime we measure---centering the distribution cannot
change \emph{how many} sequences fall outside the region, only \emph{which} ones. What
the Corrector removes is the length discrimination and the one-sidedness, and that is
the basis on which we evaluate it (\S\ref{sec:eval}). It is also scoped: under strictly
synchronous on-policy training, importance weights are unnecessary and so is the
Corrector. It earns its place in the asynchronous multi-tenant setting, where importance
weights are needed to absorb staleness and are therefore contaminated by a framework
artifact that has nothing to do with staleness. Separating the two inside the importance
weight is exactly what it does.

Finally, the Corrector is what makes the hybrid backend routable. Because it restores
every sampling path to the trainer's reference, the fixed and flex paths become
interchangeable from a tenant's point of view even though one may be quantized and the
other is not. Without it, the router could not move a tenant's rollouts between paths
without silently changing the quality of that tenant's gradients.
%======================================================================
\section{Contiguous Training Scheduling}
\label{sec:sched}

\begin{figure}
    \centering
    \includegraphics[width=\linewidth]{figures/arch_figure/Tuft-requests scheduling.pdf}
    \caption{Contiguous training scheduling. (a) Three tenants, each a sampling phase
    $S_i$ followed by a training burst $T_i$. (b) Naive scheduling lets the bursts
    collide, forcing one to wait and leaving the training lane idle between them. (c)
    Delaying the sampling of later tenants shifts their bursts so training packs
    back-to-back with no collision, opening one long concentrated sampling window.}
    \label{fig:reqsched}
    \Description{}
\end{figure}
TuFT's scheduler delays each tenant's sampling so that training bursts line up
back-to-back on a single lane. Compared with naive scheduling, which lets
bursts collide and leaves idle gaps (Fig.~\ref{fig:reqsched}(b)), our
\emph{contiguous training} scheduling (Fig.~\ref{fig:reqsched}(c)) packs
training into few contiguous blocks, opening a long concentrated sampling
window between them and minimizing switching. It has two independent control
loops that couple only through the measured training-burst duration $d_i$: a
\emph{fast loop} that assigns each tenant a one-shot sampling delay $\phi_i$
(\S\ref{sec:sched:fast}), and a \emph{slow loop} that re-composes
Fixed$\leftrightarrow$Flex groups to set how much of the cluster is training-capable
(\S\ref{sec:sched:slow}).

\subsection{Loop discovery}
\label{sec:sched:discovery}
For each tenant $i$, the scheduler infers period $P_i$, training-burst duration
$d_i$, natural release offset $r_i$, and the effective staleness window $W_i$
(implicit from which adapter versions the tenant consumes in \textsf{update};
\S\ref{sec:service-model}). Estimation runs from the raw signals of the two
API calls only---submission timestamps, per-step duration, sample volume, and
version-advance events---and yields scheduling-invariant \emph{demand}
(\S\ref{sec:rl-loop-char}). $d_i$ is always the value measured under the
\emph{current} Fixed/Flex composition; we assume no analytic relation between the
number of training-capable groups and how fast a burst executes. Drift
or tenant arrival/departure triggers a re-fit and a fast-loop replan.

\subsection{Fast loop: contiguous training packing}
\label{sec:sched:fast}

Training in TuFT is a single logical task (\S\ref{sec:backend}); the fast loop
decides where to place each tenant's training burst on this single lane so that
bursts are non-overlapping (no queueing) and packed contiguously (a long
concentrated sampling window between blocks; few switches). Its sole lever is
the sampling delay $\phi_i$: shifting a tenant's sampling later shifts its
entire loop later, which shifts when its training becomes ready.

\begin{algorithm}[t]
\small
\begin{algorithmic}[1]
    \LComment{On the scheduler, at each fast-loop replan}

    \Function{FastLoop}{$\{(P_i, d_i, r_i, W_i)\}$}
        \LComment{Step 1: Order tenants by readiness}
        \State {\em Order} $\leftarrow$ tenants sorted by ready time $r_i$ ascending
        \State {\em LaneFree} $\leftarrow -\infty$ \Comment{single training lane}

        \LComment{Step 2: Tuck each training burst against the previous one}
        \ForAll{tenant $i$ in {\em Order}}
            \State $t_i \leftarrow \max(r_i,\ \mathit{LaneFree})$ \Comment{first legal start}
            \State $\phi_i \leftarrow t_i - r_i$ \Comment{one-shot sampling delay}

            \LComment{Step 3: Reject if the delay exceeds the staleness window}
            \If{$\phi_i > W_i$}
                \State \Return \textsc{Infeasible} \Comment{lane too tight $\to$ slow loop}
            \EndIf

            \State {\em LaneFree} $\leftarrow t_i + d_i$
        \EndFor

        \LComment{Step 4: Emit the delay/start assignment}
        \State \Return $\{(\phi_i, t_i)\}$
    \EndFunction

\end{algorithmic}
\caption{Fast loop: sampling-delay assignment.}
\label{alg:space-fast}
\end{algorithm}

The fast loop (Alg.~\ref{alg:space-fast}) is release-ordered greedy packing on a single
training lane: tenants are considered in order of when their training becomes ready, and
each is given the earliest legal start, so it runs either at its natural release
(freshest) or tucked directly against the previous burst, whichever comes later. This
makes each tenant's own delay as small as it can be given that order, which is what keeps
the schedule close to the freshest feasible one; the deadline enters not as a sort key but
as the feasibility test. When some $\phi_i > W_i$---the lane cannot absorb the tenant
within its staleness window---the fast loop signals the slow loop to raise capacity rather
than violating freshness.

\subsection{Sampling-side scheduling}
\label{sec:sched:affix}

\begin{figure}
    \centering
    \includegraphics[width=\linewidth]{figures/arch_figure/Tuft-sampling requests rooter.pdf}
    \caption{The three-level sampling request pipeline: initial-adapter merging,
    request reordering, and routing whole adapter groups to fixed and flex sampling
    groups.}
    \label{fig:samplerouter}
    \Description{}
\end{figure}

The fast loop decides \emph{when} each tenant's loop should run; the sampling side has to
realize that decision across a pool of vLLM groups while keeping every group busy. Two
things make this non-trivial: multi-LoRA batching is efficient only when a step is
adapter-coherent, and a flex group's capacity is not permanent---it will be reclaimed for
training at a time the fast loop has already fixed.

Applying the delay is the easy part, and it is not one of the pipeline stages. A delay
$\phi_i$ is a \emph{relative} shift of tenant $i$'s entire loop. Because sampling and
training are chained ($\text{sample}\!\to\!\text{train}\!\to\!\text{sample}$), holding
$i$'s next request once by $\phi_i$ shifts every later iteration by the same amount; the
hold is then cleared, and no per-step throttling is needed. The delay is therefore a
one-shot gate in front of the pipeline, after which requests pass through three stages of
decreasing granularity (Fig.~\ref{fig:samplerouter}).

\noindent\textbf{L1 --- Initial-adapter merging.} A LoRA adapter is created with a
zero-initialized second factor, so at version $0$ it is exactly the identity map: a tenant
that has not yet applied an update is sampling from the frozen base. All such adapters are
therefore interchangeable, and \sysname merges them into one pseudo-adapter $A_0$ under the
condition $v_i = 0$---in Fig.~\ref{fig:samplerouter}, $R_2$ and $R_4$ become $R_0$
requests. This costs nothing and is pure gain: a step that would have carried three
distinct adapters carries one, and the $A_0$ path can bypass the LoRA kernels entirely. The
merge is transient by construction, since a tenant leaves $A_0$ as soon as it commits its
first update, and it applies only to adapters provably equal to the base, not to
warm-started ones.

\noindent\textbf{L2 --- Request reordering.} The merged stream is reordered so that
requests sharing an adapter key become contiguous, turning an interleaved arrival order
into a few long same-adapter runs. Ordering is stable within a key, so no request is
starved relative to another for the same adapter. This is what makes the downstream steps
adapter-coherent, and it is done deliberately \emph{before} routing rather than inside each
group: a group that receives an already-coherent run needs no further grouping of its own.

\noindent\textbf{L3 --- Routing with adapter groups.} The unit of routing is the
same-adapter run produced by L2, not the individual request. Keeping a run intact is what
preserves adapter-cache and prefix locality, so \sysname splits one only when no single
group can absorb it. An adapter is otherwise pinned to the group that last served it, and
re-pinning happens only on tenant arrival or departure, or on persistent imbalance. Note
that $A_0$ is routed as one group like any other rather than sprayed across the pool: its
requests are cheap precisely because they share a single identity, and splitting them would
give that back.

Which groups are eligible is where fixed and flex differ. Give each sampling group a
\emph{horizon}: for a fixed group it is unbounded, because a fixed group never changes
role; for a flex group it is the time remaining until its next scheduled flip. A run is
admitted to a flex group only if its estimated completion fits inside that horizon---the
red cross in Fig.~\ref{fig:samplerouter} marks a run rejected because there is no free time
left before the flip. The estimate uses a conservative upper bound on response length,
since admitting one over-long request delays a flip whereas rejecting one merely costs a
little utilization.

That single rule produces two behaviours we would otherwise have to implement separately.
Long or unpredictable work---multi-turn agentic rollouts, whose duration includes external
tool latency---never lands on a flex group, because only an unbounded horizon can hold it;
short work does, which is why the short requests $R_s$ in Fig.~\ref{fig:samplerouter} are
the ones routed to the flex group. And because a flex group's horizon shrinks monotonically
as its flip approaches, the admissible set shrinks with it: the group quiesces by itself and
is nearly empty when the flip arrives. No drain command is ever issued, and the KV cache
the flip discards is mostly empty---which is what makes the flip of \S\ref{sec:zerocopy}
cheap in practice and not merely in principle.

% (Stale pseudocode for the sampling-side scheduler removed; superseded by the L1/L2/L3 prose above.)

\subsection{Slow loop: Fixed$\leftrightarrow$Flex composition}
\label{sec:sched:slow}
Training in \sysname is always a single logical task: one tenant's update runs at a
time. When more than one Flex group is in the training role they are \emph{data
parallel over the batches of that one task}---they enlarge the effective batch and raise
training throughput, they do not run different tenants concurrently, and they do not
divide a tenant's burst duration $d_i$ among themselves. The slow loop therefore does
not tune a degree of parallelism on behalf of any tenant; it decides how much of the
cluster is training-capable at all.

The decision is a hysteresis on smoothed training pressure. Sustained high pressure
converts a Fixed group into a Flex group, buying training throughput at the cost of a
sampling group that would never have dropped its KV cache; sustained low pressure
converts a Flex group back into a Fixed group, buying warm-KV sampling capacity---and,
if the deployment enables it, quantized sampling, which is available only off the
zero-copy path (\S\ref{sec:zerocopy}). At least one group always remains
training-capable, so a tenant's \textsf{update} never waits for a conversion to finish.

Hysteresis is not a stylistic choice here. A Fixed$\leftrightarrow$Flex conversion
reloads the base weights and may change their numeric format, and it rebuilds the
serving runtime that a flip is careful to preserve; measured on 0.6B it costs $0.95$\,s
in one direction and $3.8$\,s in the other, two orders of magnitude more than the flip
of Table~\ref{tab:flipcost}. The two loops are separated precisely because their
transitions differ by that much: the fast loop operates on the cheap transition and may
act every period, while the slow loop operates on the expensive one and must amortize it
over many. The conversion threshold follows directly---a group is converted only if the
expected gain over its minimum dwell time exceeds the conversion cost.

The loops couple in exactly one place. Changing the composition changes how fast a
training burst executes, so after any conversion we let discovery observe a few
iterations, re-estimate $d_i$ under the new composition, and rerun the fast loop with the
new value. We never predict $d_i$ analytically from the number of groups; it is always
the measured value under the composition currently in force.

%======================================================================

\section{Implementation}
\label{sec:impl}

\sysname is built on PyTorch 2.11, vLLM 0.24 and Ray 2.50. This section covers the
service surface first, then the parts specific to this paper: the flex backend, the
Corrector, and adapter management.

\subsection{\sysname service architecture}
\label{sec:impl:service}

\paragraph{Tinker-compatible API.} \sysname reuses the Tinker request and response
types directly, so an unmodified Tinker client can drive it without recompilation. The
RL loop uses only a handful of the exposed HTTP routes: one call registers a tenant's
LoRA adapter, \textsf{forward\_backward} and \textsf{optim\_step} carry out training, a
third publishes the updated adapter as a new sampler-visible version, and
\textsf{asample} draws rollouts against it. An OpenAI-compatible router additionally
lets an adapter be addressed as a model name, so off-the-shelf agent frameworks can
drive a tenant's policy during rollout without knowing about \sysname at all. The
tenant never names a GPU, a backend or a schedule (\S\ref{sec:service-model}).

\paragraph{FutureStore.} Every mutating call returns immediately with a request
identifier; the client retrieves the result by polling. Internally a
\textsf{FutureStore} owns the lifecycle---enqueue, execute, record, expire---and
persists records to Redis. Two ordering guarantees are layered on top. Each training
run holds an execution lock and a monotone sequence counter, so a tenant's own
operations can never be reordered or interleaved, and a stale or out-of-order
submission is rejected rather than silently applied. Each sampling session has its own
sequence executor for history recording. Ordering is therefore per tenant by
construction, which is what allows the scheduler to move tenants relative to one
another freely.

\paragraph{Tenancy.} An API key maps to a user; ownership is checked on every session,
run, future, checkpoint and sampler, and cross-tenant access is possible only through
an explicit publish flag on a checkpoint. All tenants of a base model share the same
resident base weights and only differ in their LoRA adapter, so admitting a tenant
costs an adapter and its optimizer state rather than a model copy---the property the
whole system is built on.

\paragraph{Execution backends.} \sysname is engine-pluggable. On the training side we
provide three interchangeable backends: HuggingFace with PEFT for the single-actor
path, FSDP2 for sharded training, and the flex trainer of \S\ref{sec:impl:flex} for
zero-copy operation. On the sampling side we currently use vLLM, wrapped so that
prefix caching and chunked prefill are enabled and returned log-probabilities are the
ones actually sampled from. All backends are Ray actors, which gives process isolation
between trainer and sampler on the same node and a straightforward path to multi-node
placement. Which backend is used is a deployment choice; the rest of the paper focuses
on the flex trainer paired with vLLM, since that is the pairing that enables the
zero-copy switch.

\subsection{Flex backend}
\label{sec:impl:flex}

The flex backend is a mode-transition layer that sits underneath both engines and is
independent of either. Its central invariant is that at any instant exactly one copy of
the base weights exists in device memory, held by the trainer and referenced by the
sampler through an alias.

\paragraph{Trainer.} PyTorch native tensor parallelism with the fused, rank-major
layout of \S\ref{sec:zerocopy}, so no conversion is needed at alias time. Base
parameters are frozen; LoRA is attached as wrapper containers on the fused linear
layers, each holding a registry of adapters with one active at a time; the optimizer
covers only LoRA parameters. At $\text{TP}=1$ the parallelization is skipped so a single
GPU can run the whole loop.

\paragraph{Alias injection.} The sampler starts with a dummy weight load and vLLM's
pluggable virtual-memory allocator enabled, so weights and KV cache draw from a pool
that supports unmap and remap. Trainer shards are exported as IPC descriptors and each
sampler worker rebinds its parameters to the producer's pages. Injection runs from a
custom worker class that overrides the warm-up entry point, placing it after KV cache
initialization and before graph capture; capturing the graph after the alias is what
keeps the alias valid across every subsequent flip. Displaced dummy allocations are
re-tagged so the allocator continues to track their virtual addresses.

\paragraph{Sleep and wake.} vLLM's native sleep offloads weights to host memory, which
is incompatible with an alias, so \sysname implements its own. Entering training unmaps
the physical pages behind the KV cache and the re-tagged dummy weights, leaving the
graph pool resident; returning to sampling re-maps the same handles. On 4B at
$\text{TP}=2$ this releases 39.1\,GB per rank in 96\,ms and restores it in 29\,ms; on
32B at $\text{TP}=4$, releasing 12\,GB costs 144\,ms out and 732\,ms back.

\paragraph{A LoRA-cache correctness bug.} Our first flex implementation trained
measurably worse than a HuggingFace-plus-vLLM baseline, and the cause was neither the
trainer nor the switch. Under the fast sleep/wake path vLLM could serve a \emph{stale}
LoRA adapter from its internal cache: changing the numeric adapter identifier was not
enough to force a refresh while the adapter name stayed fixed. The sampler was
therefore generating from an older policy than the trainer scored, and importance
ratios diverged---log absolute mean $0.79$ with ratios up to $55$. Making the internal
adapter name unique per exported snapshot while keeping the tenant-visible identifier
stable fixed it, restoring the ratio to $0.02$ with a maximum of $1.8$; on a 50-step
Countdown RL run the reward rose from $0.10$ to $0.37$, against $0.28$ for the
baseline. We report it because it is invisible to any test of switching latency or
memory and only shows up in convergence.

\paragraph{Router to the fixed backend.} Because quantization would break the alias,
the fixed sampling backend is a separate vLLM runtime that loads its own base copy and
may enable FP8 or other schemes independently. A router owns the transition: moving to
fixed first shrinks the flex group to its training footprint, then wakes the fixed
runtime; moving back puts the fixed runtime to sleep before flex sampling expands
again. Known adapters are replayed into whichever runtime becomes active. This
transition is the slow loop of \S\ref{sec:sched:slow}; on 0.6B it costs $0.95$\,s in
one direction and $3.8$\,s in the other, dominated by waking the fixed runtime.

Two implementation constraints follow directly from the alias: vLLM's custom all-reduce
fails to capture a graph over aliased weights, so both the flex path and the sampling
baseline run with it disabled; and the vocabulary-sharded embedding and output
projection are not aliased because vLLM's shard does not correspond to the trainer's
full-vocabulary tensor.

\subsection{Corrector}
\label{sec:impl:corrector}

The Corrector is a five-million-parameter model dominated by its token embedding, in
two variants: an MLP scoring each token independently and a two-layer transformer that
adds context. Both take the sampling-side features of \S\ref{sec:mismatch-pred} and
emit a per-token residual whose sequence mean is the applied correction. The output
head is zero-initialized so an untrained Corrector is bit-identical to no correction,
which makes the ablation exact.

Calibration data is collected from synchronous rollouts, where staleness is zero by
construction and the observed gap is entirely framework-induced. Splits are taken by
weight version and by tenant rather than at random, since a random split would place
tokens from the same policy version on both sides and inflate the result. Training is
inexpensive---minutes on CPU for the MLP variant---and is performed once per base model
per sampling path. At serving time the Corrector runs on the sampler, consumes features
that are already available, and requires no extra forward pass; its output is detached
so no gradient can flow back into it.

Two numerical prerequisites make the mismatch measurable in the first place. vLLM does
not apply temperature scaling when computing prompt log-probabilities, so we patch the
model runner to divide logits by the sampling temperature and require the trainer to
apply the same scaling before its log-softmax. And we disable top-$p$ and top-$k$
during rollout, since under truncated sampling the reported log-probability is no
longer the log-probability of the distribution actually sampled from. Without both, the
reported gap is dominated by a deterministic temperature transform rather than by any
real framework difference.

\subsection{Adapter management}
LoRA adapters are small enough that the storage path is not the bottleneck. Measured
on our testbed, Linux page cache, local disk and shared storage (NAS) all provide
sufficient bandwidth for load, publish and synchronization; the total adapter-related
overhead is $3\%\!\sim\!6\%$ of an RL iteration on NAS. A new adapter version reaches
the sampler through disk rather than through a weight-transfer path: the trainer
exports a standard PEFT adapter, the platform records it as a versioned checkpoint
under a stable handle, and the sampler loads it as a LoRA request. Adapter slots on the
training side are pre-allocated per rank and reset on release---reinitializing the
low-rank factors, dropping the optimizer state and clearing pending gradients---so that
admitting a tenant does not allocate and departure does not fragment.

%======================================================================
\section{Evaluation\ \todo{(mock now)}}
\label{sec:eval}

Our evaluation answers three questions:

\begin{itemize}
    \item Under a fixed GPU budget, does \sysname sustain higher RL fine-tuning
    throughput than existing designs, and how far does it scale in the number of
    concurrent tenants?
    \item Does it hold the three per-tenant service SLOs---response time, fidelity and
    freshness---while doing so?
    \item How much does each mechanism contribute?
\end{itemize}

\subsection{Experimental setup}

\paragraph{Testbed.} A single node with 8 A100-80GB GPUs, NVLink-connected. All systems
run the same PyTorch, vLLM and CUDA versions, and---because aliased weights are
incompatible with vLLM's custom all-reduce (\S\ref{sec:impl})---all systems run with
custom all-reduce disabled, so the sampling baseline is measured under exactly the
configuration \sysname uses.

\paragraph{Workloads.} We use Qwen3-4B for the main experiments and Qwen3-32B for the
scale experiments, with up to 32 concurrent tenants sharing one frozen base and one LoRA
adapter each. Tenant tasks mix single-turn reasoning (GSM8K, MATH) with multi-turn
agentic workloads that issue tool calls, giving heterogeneous response lengths,
rollout-to-training ratios and loop periods. All systems read the same tenant
configuration file, use GRPO with identical hyper-parameters, and are held to the same
staleness budget of three policy versions: a rollout is admitted for training only while
its version lag is within the bound, and a tenant that runs further ahead stops
sampling. We treat this bound as a property of tenant behaviour rather than a platform
knob, so it applies identically to every system.

\paragraph{Baselines.} We compare against four designs, each implemented as a policy
inside our own runtime so that the model, kernels, engine versions and RL algorithm are
identical and only the resource model differs. Table~\ref{tab:baselines} summarizes them.

\begin{table}[t]
\centering
\footnotesize
\setlength{\tabcolsep}{4pt}
\caption{Baseline configurations. All four share the same base model, RL algorithm and
staleness budget; they differ only in how sampling and training capacity is organized.}
\label{tab:baselines}
\begin{tabular}{lcccc}
\toprule
 & Engines & Base & Boundary & Shared \\
 &         & copies &        & sampling \\
\midrule
Serial-Async     & 2 & 2 & static & no  \\
Unified-Engine   & 1 & 1 & shared & yes \\
Colocate-2Copies & 2 & 2 & shared & yes \\
Static-Disagg    & 2 & 2 & static & yes \\
\midrule
\sysname         & 2 & \textbf{1} & \textbf{elastic} & yes \\
\bottomrule
\end{tabular}
\end{table}

\emph{Serial-Async}~\cite{verlhybridflow} is the status quo: each tenant runs its own
asynchronous RL loop on a statically partitioned cluster. Because a single node cannot
host 32 independent jobs, tenants time-share the cluster, each owning all GPUs for $k$
RL iterations before yielding; the frozen base stays resident and only LoRA weights and
optimizer state are swapped, which is more favourable than tearing down and reloading
each job. Only one adapter is active in the sampler at a time, so there is no
cross-tenant overlap of any kind. We sweep the train:sample split and $k$ and report the
best-throughput configuration.

\emph{Unified-Engine} removes the second engine: a single tensor-parallel trainer both
samples and trains, alternating between the two modes when a training queue fills.
Mode changes require no resharding and no re-initialization, and the log-probabilities
served to tenants are by construction identical to those the trainer computes, so this
baseline has \emph{zero} framework mismatch. Its cost is sampling throughput, since a
training runtime is not an inference engine. We enable batched decode, a preallocated KV
cache, prefill reuse and adapter-homogeneous batching so that this comparison is not
against a strawman generator. We adopt the co-serving principle
of FlexLLM~\cite{oliaro2025flexllmtokenlevelcoservingllm} rather than its mechanism:
FlexLLM fuses fine-tuning tokens into an inference batch inside an inference engine, and
that direction is not reproducible here because vLLM has no backward pass, nor is its
premise satisfied in RL, where a tenant's training work does not exist until its own
rollouts finish.

\emph{Colocate-2Copies} keeps both engines co-resident on the same GPUs, each holding its
own copy of the base weights, and alternates with the same queue-triggered
policy~\cite{verlhybridflow}. Switching is inexpensive because both runtimes stay
resident and the KV cache is retained across a training pause, but the duplicated base
leaves substantially less memory for KV, capping sampling batch size. This baseline
isolates the value of \sysname's zero-copy weight sharing, and it retains the framework
mismatch inherent in running two engines.

\emph{Static-Disagg}~\cite{yu2026marlaasmultitenantasynchronousreinforcement} splits the
cluster once at deployment into a minimal training group and a sampling group that
receives all remaining GPUs, and never changes it. The sampler keeps every tenant's
adapter resident and batches across tenants, admission is throttled by an estimate of
aggregate KV footprint, and training is serialized through a global FIFO. We sweep its
split independently and report its best configuration. We also grant it the same
staleness budget of three versions that \sysname exploits; the original design is
strictly on-policy and would therefore be strictly slower.

For both static baselines the split is not chosen by hand: we enumerate every feasible
train:sample partition and report each baseline at its own best point, whereas \sysname
receives an arbitrary initial composition and is not tuned per workload. We verify that
Serial-Async at a single tenant matches a stock verl deployment to within a few percent,
so that no multi-tenant gain is attributable to a slow baseline.

\subsection{Overall performance}

\paragraph{Throughput under a fixed GPU budget.} We fix the cluster at 8 GPUs, run the
32-tenant workload for a fixed wall-clock window after warm-up, and report aggregate
training steps per hour together with accelerator utilization and idle fraction.
\todo{\sysname reaches $\mathbf{[\cdot]}$ steps/hour against $\mathbf{[\cdot]}$ for the
strongest baseline, a $\mathbf{[\cdot]}\times$ improvement, and raises utilization from
$\mathbf{[\cdot]}\%$ to $\mathbf{[\cdot]}\%$.} The same experiment is repeated on
Qwen3-32B, where the feasible static partitions collapse to a single point: at that
scale a static system has no tuning freedom left at all, while \sysname can still move
its boundary.

\paragraph{Scaling with tenant count.} Sweeping from 1 to 32 tenants exposes the
capacity knee. Prior multi-tenant designs degrade beyond roughly four concurrent
tenants~\cite{yu2026marlaasmultitenantasynchronousreinforcement};
\todo{\sysname moves the knee from ${\sim}\mathbf{X}$ to ${\sim}\mathbf{Y}$, and spatial
isolation cannot be provisioned past $\mathbf{[\cdot]}$ tenants on this testbed,} since
each isolated job needs its own base copy, so the Serial-Async curve is truncated there
and continues as its time-shared variant.

\begin{figure}[t]
    \centering
    \begin{minipage}[t]{0.48\linewidth}
        \centering
        \includegraphics[width=\linewidth]{figures/exp_figure/eval_throughput_gpu_budget.pdf}
        \caption{Throughput under fixed GPU budgets.}
        \Description{Grouped bar chart comparing training throughput of Serial-Async,
        Unified-Engine, Colocate-2Copies, Static-Disagg and TuFT under 4-GPU and 8-GPU
        budgets.}
        \label{fig:eval_throughput_gpu_budget}
    \end{minipage}
    \hfill
    \begin{minipage}[t]{0.48\linewidth}
        \centering
        \includegraphics[width=\linewidth]{figures/exp_figure/eval_scaling_tenants.pdf}
        \caption{Scaling with concurrent tenants.}
        \Description{Line chart showing throughput versus tenant count for the baseline
        policies and TuFT, with TuFT sustaining higher throughput at larger tenant
        counts.}
        \label{fig:eval_scaling_tenants}
    \end{minipage}
\end{figure}

\subsection{Per-tenant serving SLOs}

Averages hide exactly the failures these SLOs exist to prevent, so every quantity in
this subsection is reported as a distribution: median and 95th percentile, with the full
CDF where it is informative, plus the fraction of requests that meet the target.

\subsubsection{Response time}
We separate the two latencies a tenant observes. \emph{Time-to-first-step} measures the
delay between a tenant's training submission and the moment its update begins executing;
it is the quantity that time-shared designs sacrifice, since a tenant there waits
roughly a full rotation for its turn. \emph{Sampling response time} measures time to
first token and end-to-end rollout latency, and it is the quantity that switch-based
designs sacrifice, since their sampling stops entirely whenever training runs.

\subsubsection{Fidelity}
We report the distribution of the absolute difference between the log-probability served
to a tenant and the one its trainer computes for the same token, together with the
fraction of sequences whose accumulated sequence-level bias falls within a target
$\epsilon$. Both \sysname sampling paths are reported separately, since the fixed path
may be quantized. Unified-Engine appears here as the zero-mismatch reference point: it
achieves perfect fidelity by construction and pays for it in throughput, which makes the
trade-off this SLO captures explicit rather than rhetorical.

\begin{figure}[t]
    \centering
    \begin{minipage}[t]{0.46\linewidth}
        \centering
        \includegraphics[width=\linewidth]{figures/exp_figure/eval_response_time.pdf}
        \caption{Response time.}
        \Description{Grouped box plot comparing training and sampling latency
        distributions across all baseline policies and TuFT.}
        \label{fig:eval_response_time}
    \end{minipage}
    \hfill
    \begin{minipage}[t]{0.46\linewidth}
        \centering
        \includegraphics[width=\linewidth]{figures/exp_figure/eval_fidelity_mean_logprob_mismatch.pdf}
        \caption{Log-probability mismatch.}
        \Description{Point chart with error bars showing mean and 95th percentile
        absolute log-probability mismatch for each system on the same token set.}
        \label{fig:eval_fidelity_mean_logprob_mismatch}
    \end{minipage}
\end{figure}

\subsubsection{Freshness}
We report the distribution of version lag at the moment each rollout is consumed for
training, the fraction of tenants that ever stall against the staleness bound, and the
GPU-time wasted by such stalls. Because the bound is enforced identically everywhere,
this table also serves as the audit that no system won by quietly relaxing it.

\begin{figure}[t]
    \centering
    \includegraphics[width=\linewidth]{figures/exp_figure/eval_freshness.pdf}
    \caption{Average parameter update version observed by each tenant. Lower spread
    indicates that tenants progress more uniformly under the same staleness bound.}
    \Description{Line chart comparing the average parameter update version observed by
    each tenant for the baseline policies and TuFT.}
    \label{fig:eval_freshness}
\end{figure}

\subsection{Mechanism microbenchmarks}

\subsubsection{Backends performance: working and switching}

\begin{figure}[t]
    \centering
    \includegraphics[width=\linewidth]{figures/exp_figure/eval_backend_breakdown.pdf}
    \caption{Backend working and switching breakdown. Bars decompose each iteration
    segment into training compute, sampling decode, queue waiting and switching cost.}
    \Description{Horizontal stacked bar chart decomposing backend iteration time into
    training compute, sampling decode, queue wait and switching overhead.}
    \label{fig:eval_backend_breakdown}
\end{figure}


\subsubsection{Corrector performance}

\begin{figure}[t]
    \centering
    \includegraphics[width=\linewidth]{figures/exp_figure/eval_corrector_performance.pdf}
    \caption{Corrector performance over generations. Curves report the average
    percentage reduction in sequence-level bias relative to uncorrected two-engine
    log-probabilities.}
    \Description{Line chart showing average sequence bias reduction percentage over
    generations for token correction, sequence correction and adaptive correction.}
    \label{fig:eval_corrector_performance}
\end{figure}

\subsubsection{GPU utilization with scheduling}

\begin{figure}[t]
    \centering
    \includegraphics[width=\linewidth]{figures/exp_figure/eval_gpu_utilization.pdf}
    \caption{GPU utilization over time for the baseline policies and \sysname.}
    \Description{Time-series line chart comparing GPU utilization over time for each
    baseline policy and TuFT.}
    \label{fig:eval_gpu_utilization}
\end{figure}

\subsection{Ablation}

The throughput--quality plane in Fig.~\ref{fig:eval_ablation_quality_throughput} pins
down where each mechanism contributes. Starting from a static disaggregation baseline,
we add \sysname's mechanisms one at a time and ask how the point moves.

\emph{TuFT-noSched} keeps the hybrid backends but removes the scheduler: each flex
group flips as soon as any tenant has a full training batch, sampling requests are
routed without horizon-aware admission (\S\ref{sec:sched:affix}), and there is no delay
lever. This variant isolates the raw value of moving the boundary---it is what static
disaggregation would look like if its split were free to change.
\emph{$+$Fast loop} enables one-shot delay and release-ordered packing on a single
training lane (\S\ref{sec:sched:fast}); training bursts stop colliding and switches
become rare, though sampling routing is still adapter-blind.
\emph{$+$Sampling routing} turns on adapter-group routing with horizon-fit admission
(\S\ref{sec:sched:affix}); flex groups quiesce before their flip and same-adapter runs
stay contiguous.
\emph{$+$Slow loop} enables Fixed$\leftrightarrow$Flex re-composition
(\S\ref{sec:sched:slow}), letting the training-capable share of the cluster track
sustained load.
\emph{$+$Corrector} completes full \sysname by restoring log-probability fidelity on
both sampling paths (\S\ref{sec:mismatch-pred}); it barely moves the throughput axis
but carries a large step along the quality axis, since it removes the systematic length
bias in the importance weight.

Each purple point in Fig.~\ref{fig:eval_ablation_quality_throughput} therefore
corresponds to one mechanism. The movement decomposes cleanly: fast and slow scheduling
improve throughput at essentially unchanged quality; the Corrector improves quality at
essentially unchanged throughput; zero-copy switching (already inside every variant) is
the only mechanism that moves both axes at once. No single mechanism gives the whole
gain, which is why the ablation is a Pareto frontier rather than a single bar.

\begin{figure}[t]
    \centering
    \includegraphics[width=\linewidth]{figures/exp_figure/eval_ablation_quality_throughput.pdf}
    \caption{Throughput--quality trade-off for baselines and \sysname variants. The
    purple points show incremental TuFT mechanisms, while baseline points provide the
    comparison frontier.}
    \Description{Scatter plot of throughput versus training quality for baseline
    policies and TuFT variants, with TuFT variants shown as purple points.}
    \label{fig:eval_ablation_quality_throughput}
\end{figure}


%======================================================================
